\documentclass[11pt]{article}
\usepackage[T1]{fontenc}
\usepackage{lmodern,amsmath,amssymb,amsthm,mathtools,microtype,xcolor}
\usepackage[margin=1in]{geometry}
\usepackage[colorlinks=true,linkcolor=blue!45!black,citecolor=blue!45!black,urlcolor=blue!45!black]{hyperref}
\newtheorem{theorem}{Theorem}[section]
\newtheorem{lemma}[theorem]{Lemma}
\newtheorem{corollary}[theorem]{Corollary}
\newtheorem{proposition}[theorem]{Proposition}
\theoremstyle{definition}
\newtheorem{sourceinput}[theorem]{Input}
\theoremstyle{remark}
\newtheorem{remark}[theorem]{Remark}
\newtheorem{example}[theorem]{Example}
\newcommand{\kk}{\overline{\mathbb Q}}
\newcommand{\Qa}{\overline{\mathbb Q}}
\newcommand{\Kexp}{K_{\mathrm{exp}}}
\newcommand{\VM}{\operatorname{VM}}
\newcommand{\Sone}{\mathbb S^1}
\newcommand{\Ex}{\mathbb E}
\newcommand{\trdeg}{\operatorname{trdeg}}
\setlength{\emergencystretch}{3em}
\hypersetup{pdftitle={Cross-field Consequences of Bessel-Value Independence},pdfauthor={Keio}}
\title{Cross-field Consequences of Bessel-Value Independence:\\
Circular Mixtures, Random-Time Quantum Walks,\\and Robin Spectra}
\author{Keio}
\date{22 October 2026\\\small Research transfer draft}
\begin{document}
\maketitle
\begin{abstract}
We derive applications in statistics, spectral graph theory, and elliptic
spectral theory from explicitly stated arithmetic inputs of the unified
Euler--Bessel manuscript. For finite von Mises mixtures with algebraic
weights, positive algebraic concentrations, and algebraic unit directions,
the first $r$ circular moments identify the entire mixing measure when
each concentration supports at most $r$ directions. This bound is sharp
and independent of the number of concentrations; its inverse is nevertheless
discontinuous. We also exclude all-algebraic interior likelihood critical
points. For finite algebraic Hermitian Hamiltonians sampled at beta-distributed
random times of algebraic half-width, every averaged mixing eigenvalue
other than $1$ is transcendental over the field of algebraic exponentials;
all principal submatrices are nonsingular. For Robin problems on balls,
integral-intersection and unit inputs exclude positive algebraic eigenvalues
for a wider ring of boundary parameters, with explicit negative exceptions
in odd dimensions. A recurrence argument separately removes accidental
angular degeneracy for every algebraic Robin parameter, allowing the
Hassannezhad--Sher Pleijel theorem to apply without a genericity restriction
in that arithmetic subclass. We distinguish these applications from the
general open problems and make no claim of established historical priority.
\end{abstract}

\paragraph{Status and dependencies.}
The arguments below are complete deductions from the stated inputs of
\cite{Unified}, a working manuscript rather than a published theorem source.
The statistical and quantum-walk results use only Input~\ref{xf:input}.
The stronger Robin results also use the ring statements in
\eqref{sp:inputs}. The algebraic-parameter nondegeneracy proof is a
classical recurrence/transcendence argument. Literature comparison is
limited to the cited primary sources and does not establish novelty.
In particular no claim is made to solve the unrestricted random-time
quantum-walk programme, the general-domain Pleijel question, or the
clamped-plate multiplicity conjecture.

\section{Arithmetic input and statistical conventions}


Write $k=\overline{\mathbb Q}$ and
$\Kexp=\kk(e^a:a\in\kk)$.
We isolate precisely the value-independence statement used here.

\begin{sourceinput}[Joint Bessel values from the unified manuscript]
\label{xf:input}\label{cf:joint}
For any nonzero algebraic numbers $t_1,\ldots,t_m$ distinct modulo sign,
the $2m$ numbers $I_0(t_j),I_1(t_j)$ are jointly algebraically
independent over $\Kexp$.
\end{sourceinput}

This is a subfamily of \cite[Theorem 3.3]{Unified}.
The cited proof uses the product $\mathrm{SL}_2^m$ differential Galois
group of the distinct Bessel scales, its separation from solvable
exponential systems, and Siegel--Shidlovskii specialization at an ordinary
algebraic point. The present note proves applications of this input;
it does not reproduce or independently certify that foundational proof.
For positive distinct algebraic $t_j$, the birational change of coordinates
gives the equivalent independence of
\begin{equation}\label{vm:input}
 I_0(t_1),C(t_1),\ldots,I_0(t_m),C(t_m),
 \qquad C(t)=I_1(t)/I_0(t).
\end{equation}
Consequently $1,C(t_1),\ldots,C(t_m)$ are linearly independent over
$\kk$, and each $C(t_j)$ is transcendental over
$\Kexp(I_0(t_1),\ldots,I_0(t_m))$.
The statistical results use this positive-scale consequence.
The quantum-walk application also uses imaginary scales.
Section~\ref{sp:section} states the two additional ring-theoretic inputs
needed for its stronger Robin spectral assertions.

Identify $\Sone$ either with complex numbers of modulus one or with
real unit vectors in $\mathbb R^2$.  For $u\in\Sone$ and $t>0$, let
$\VM(u,t)$ have density, relative to uniform probability measure on the
circle,
\[
 f_{u,t}(z)=\frac{\exp(t\operatorname{Re}(\overline u z))}{I_0(t)}.
\]
Its complex circular moments are
\begin{equation}\label{vm:moment}
 \int_{\Sone} z^n\,d\VM(u,t)(z)=u^nR_n(t),\qquad
 R_n(t)=I_n(t)/I_0(t),\quad n\geq0.
\end{equation}
This follows directly by rotating the standard integral for $I_n$.
The concentration MLE and its Bessel-ratio score are standard; see
Hornik--Gr\"un~\cite{HG}.  Ordinary finite-mixture identifiability is
already classical~\cite{FHW,Kent}.  Our question is how many exact
moments suffice under explicit arithmetic restrictions.

An \emph{arithmetic mixture} means
\begin{equation}\label{vm:mixture}
 P=\sum_{j=1}^{m}w_j\VM(u_j,t_j),\qquad
 w_j\in\kk\cap\mathbb R_{>0},\quad \sum_jw_j=1,
 \quad u_j\in\kk\cap\Sone,
 \quad t_j\in\kk\cap\mathbb R_{>0}.
\end{equation}
Repeated identical pairs $(u_j,t_j)$ are merged.  The direction condition
means that the Cartesian coordinates of $u_j$ are algebraic; it does
\emph{not} mean that an angle chosen to represent $u_j$ is algebraic.
Equality of mixtures below means equality of their discrete mixing
measures, up to permutation, with identical pairs merged.

\section{A sharp moment bound independent of the number of concentrations}

\begin{lemma}[Separation of concentration groups]\label{vm:separate}
Let two arithmetic mixtures have equal $n$th complex circular moment,
where $n\geq1$.  For each positive concentration $t$ occurring in either
mixture, their grouped directional moments
\[
 a_{t,n}:=\sum_{j:t_j=t}w_ju_j^n
\]
are equal.  An absent group has all its moments zero.
\end{lemma}
\begin{proof}
The Bessel recurrence gives
\[
 R_{n+1}(t)=R_{n-1}(t)-\frac{2n}{t}R_n(t),\qquad R_0=1,\quad R_1=C(t).
\]
Thus $R_n(t)=A_n(t)+B_n(t)C(t)$, with
$A_n,B_n\in\mathbb Q[t^{-1}]$, $B_0=0$, and $B_1=1$.
For $t>0$, $B_n(t)$ has sign $(-1)^{n-1}$ for every $n\geq1$.
Indeed $B_2=-2/t<0$, and in the recurrence for $B_{n+1}$ the two
terms $B_{n-1}$ and $-(2n/t)B_n$ have the same sign.  In particular,
$B_n(t)\ne0$.

Write $\Delta a_{t,n}$ for the difference of the two grouped moments.
The assumed equality is
\[
 0=\sum_t A_n(t)\Delta a_{t,n}
     +\sum_t B_n(t)\Delta a_{t,n}C(t).
\]
All coefficients are algebraic.  Linear independence of $1$ and the
distinct $C(t)$, from \eqref{vm:input}, gives
$B_n(t)\Delta a_{t,n}=0$ for each $t$, proving the assertion.
\end{proof}

\begin{lemma}[Positive circular atomic measures with unknown mass]
\label{vm:toeplitz}
Two finite positive measures on $\Sone$, each supported on at most $r$
points, are equal if their moments of orders $1,\ldots,r$ are equal.
Their total masses need not be specified in advance.
\end{lemma}
\begin{proof}
Write $a_n=\int u^n\,d\nu(u)$ and $a_{-n}=\overline{a_n}$.
The Toeplitz matrix
\[
 T=(a_{j-i})_{0\leq i,j\leq r}=a_0I_{r+1}+H
\]
is positive semidefinite and has rank at most $r$, because it is a sum
of at most $r$ positive rank-one matrices.  The Hermitian matrix $H$
is determined by $a_1,\ldots,a_r$ and has zero diagonal.  Singularity
and positivity therefore force
\[
 a_0=-\lambda_{\min}(H).
\]
The two measures have the same total mass as well as moments
$-r,\ldots,r$.  Their signed difference is supported at $s\leq2r$
distinct nonzero nodes $v_1,\ldots,v_s$.  Using any $s$ consecutive
exponents in $[-r,r]$ gives a Vandermonde system with nonzero
determinant.  Every signed weight is zero.
The argument also includes the zero measure.  This is a direct form
of the classical positive Toeplitz/Vandermonde mechanism; compare
\cite{YX}.
\end{proof}

\begin{theorem}[Sharp arithmetic finite-moment identifiability]
\label{vm:sharp}
Fix an integer $r\geq1$.  Among arithmetic mixtures having at most
$r$ distinct directions at each concentration, the first $r$ complex
circular moments determine the entire mixing measure.  The total number
of distinct concentrations may be any finite number and is not assumed
known.  The bound $r$ cannot uniformly be replaced by $r-1$.
\end{theorem}
\begin{proof}
Suppose two mixtures have the same moments of orders $1,\ldots,r$.
Lemma~\ref{vm:separate} gives equality of these grouped directional
moments at each concentration in the union of their concentration
sets.  At each fixed concentration, the positive directional measures
have at most $r$ atoms.  Lemma~\ref{vm:toeplitz} makes them equal,
including their masses.  This also rules out a concentration appearing
in just one mixture.  All components and weights therefore agree.

For sharpness take concentration $t=1$, equal weights $1/r$ on the
$r$th roots of unity, and compare with the same configuration rotated
by $\zeta_{2r}=e^{\pi i/r}$.  For both mixtures the moments of orders
$1,\ldots,r-1$ vanish, whereas their $r$th moments are respectively
$R_r(1)$ and $-R_r(1)$.  These differ because $I_r(1)>0$.
For $r=1$, the assertion that zero moments cannot identify a distribution
is interpreted literally.
\end{proof}

\begin{corollary}[One complex moment]\label{vm:first}
If an arithmetic mixture has no repeated concentration, its first
complex circular moment determines its number of components, all its
concentrations, all its directions, and all its weights within this
class.  In particular this is an exact injectivity theorem for mixtures
with arbitrarily many components.
\end{corollary}
\begin{proof}
Apply Theorem~\ref{vm:sharp} with $r=1$.  More directly, linear
independence of the $C(t)$ determines the coefficient $wu$ attached
to each concentration.  Since $w>0$ and $|u|=1$, it determines
$w=|wu|$ and then $u$.
\end{proof}

\begin{remark}[What is new in the transfer]
Finite von Mises mixture identifiability from the entire density is
not a new result.  The arithmetic separation in
Lemma~\ref{vm:separate}, combined with the classical positive moment
argument, gives the sharper conclusion from only $r$ moments under the
stated restrictions.  The arithmetic input is essential to that
conclusion.  A literature search did not locate this exact formulation;
this is not a claim that a previously published open conjecture has
been resolved.
\end{remark}

\section{Why exact identification does not give stable recovery}

\begin{proposition}[Discontinuity already for one moment]
\label{vm:unstable}
On the class in Corollary~\ref{vm:first}, the inverse of the first-moment
map need not be continuous, even if the output is measured as a
probability density in total variation.
\end{proposition}
\begin{proof}
Choose positive algebraic $a<c<b$ and the common direction $u=1$.
The function $C$ is strictly increasing, since
$C'(t)=\operatorname{Var}_{\VM(1,t)}(\operatorname{Re}z)>0$.
Put
\[
 w_* = \frac{C(b)-C(c)}{C(b)-C(a)}\in(0,1),
\]
and choose rational $w_n\in(0,1)$ tending to $w_*$.  The arithmetic
mixtures
$P_n=w_n\VM(1,a)+(1-w_n)\VM(1,b)$ have first moments tending to
the first moment of $P=\VM(1,c)$.
Their densities converge uniformly to that of
$P_*=w_*\VM(1,a)+(1-w_*)\VM(1,b)$.
This limiting density differs from that of $P$: otherwise a nontrivial
linear relation among $e^{ax},e^{cx},e^{bx}$ would hold for
$-1\leq x\leq1$, contrary to linear independence of distinct
exponentials.  Hence
$\|P_n-P\|_{\mathrm{TV}}\to\|P_*-P\|_{\mathrm{TV}}>0$.
\end{proof}

Thus the theorem gives no stable algorithm for reconstruction from
rounded moments, no sample-complexity bound, and no effective separation
bound.  Such quantitative statements require further arithmetic input,
for example explicit lower bounds with degree and height restrictions.

\section{Transcendental coordinates at likelihood critical points}

For observations $x_1,\ldots,x_N\in\Sone\subset\mathbb R^2$, use
the relative-uniform density and write
\[
 \ell(w,u,t)=\sum_{l=1}^N
 \log\left(\sum_{j=1}^m
     w_j\frac{e^{t_j u_j\cdot x_l}}{I_0(t_j)}\right).
\]
An interior point means $w_j>0$, $\sum_jw_j=1$, $t_j>0$, with
$u_j\in\Sone$.  Concentrations are independently variable, and repeated
concentration values at a point are allowed.

\begin{theorem}[No all-algebraic interior critical point]
\label{vm:critical}
Suppose $N\geq1$ and every Cartesian coordinate of every observation
$x_l$ is algebraic.  Then every interior critical point of $\ell$ has
at least one transcendental coordinate among its weights,
concentrations, and Cartesian direction coordinates.
\end{theorem}
\begin{proof}
Assume all those parameter coordinates are algebraic at an interior
critical point.  Let $T$ be the finite set of its distinct concentrations,
and put
\[
 F=\Kexp\bigl(I_0(t):t\in T\bigr).
\]
Each exponent $t_j u_j\cdot x_l$ is algebraic.  Therefore every
responsibility
\[
 \rho_{lj}=
 \frac{w_j e^{t_j u_j\cdot x_l}/I_0(t_j)}
 {\sum_h w_h e^{t_h u_h\cdot x_l}/I_0(t_h)}
\]
belongs to $F$ and is strictly positive.  Differentiating in $t_j$
gives
\[
 \frac{\partial\ell}{\partial t_j}
 =\sum_l\rho_{lj}(u_j\cdot x_l)
       -C(t_j)\sum_l\rho_{lj}.
\]
The first sum belongs to $F$ and the second coefficient belongs to
$F\cap\mathbb R_{>0}$.  A zero derivative would force $C(t_j)\in F$,
contrary to \eqref{vm:input}.  In fact none of the concentration
derivatives can vanish at an all-algebraic interior point.
\end{proof}

\begin{corollary}[Penalized likelihood]\label{vm:penalty}
Let $p$ be a differentiable penalty such that all its concentration
derivatives are algebraic whenever the full parameter tuple is
algebraic.  Then Theorem~\ref{vm:critical} also holds for $\ell+p$.
In particular it holds for
\[
 p=-\psi\sum_j t_j,\qquad \psi\in\kk\cap\mathbb R_{>0}.
\]
\end{corollary}
\begin{proof}
The additional derivative belongs to $\kk\subset F$ and changes only
the first sum in the proof.  The positive coefficient of $C(t_j)$
is unchanged.
\end{proof}

The linear penalty is the form used by Ng~\cite[Section 4.3]{Ng},
including $\psi=\zeta/N$ with positive algebraic $\zeta$.
Only interior critical points are covered: boundary optima, zero
weights, and $t_j=0$ require separate treatment.  Unpenalized mixture
likelihoods can be unbounded~\cite[Theorem 3.1]{Ng}, so existence of a
global ordinary MLE is not asserted.  The conclusion is that at least
one coordinate is transcendental, not that every concentration or every
coordinate is transcendental.  It is a circular-distribution analogue
of the arithmetic perspective of Am\'endola--Drton--Sturmfels
on Gaussian mixture likelihoods~\cite{ADS}.

\section{Boundaries and further questions}

The moment theorem and the likelihood obstruction use different parts
of \eqref{vm:input}: linear independence of the ratios suffices for
moment separation, whereas the likelihood argument needs independence
of the ratios from all raw normalizing constants and algebraic
exponentials.  They do not prove a stable mixture estimator, classify
likelihood critical points, establish D-finiteness or non-D-finiteness
of the inverse Bessel-ratio function, or settle general real-parameter
finite-moment identifiability.

A concrete next question is quantitative: with bounds on degrees,
heights, component weights and concentration separation, can one give
effective lower bounds for the difference of the first $r$ moment
vectors of two distinct arithmetic mixtures?  The qualitative
independence input alone supplies no such bound, and
Proposition~\ref{vm:unstable} precludes a uniform result without
additional restrictions.


\section{Arithmetic spectra of random-time quantum walks}
\label{qmix:section}

Let $H\in M_N(\Qa)$ be Hermitian, and put
\[
 U_H(t)=e^{itH},\qquad
 M_H(t)=U_H(t)\circ\overline{U_H(t)},
\]
where $\circ$ denotes entrywise multiplication. For an integer $\nu\ge0$
and $a>0$, consider the probability density
\begin{equation}
\label{qmix:density}
 f_{a,\nu}(r)=
 \frac{\Gamma(\nu+1)}{a\sqrt\pi\,\Gamma(\nu+\tfrac12)}
 \left(1-\frac{r^2}{a^2}\right)^{\nu-1/2}
 \mathbf1_{(-a,a)}(r).
\end{equation}
The cases $\nu=0,1$ are the arcsine and semicircle laws, respectively.
Write $\widehat M_{a,\nu}(H)=\Ex[M_H(R)]$ for this sampling law.
Throughout this section $\Kexp=\Qa(e^\alpha:\alpha\in\Qa)$.

\begin{theorem}[Arithmetic spectrum under beta sampling]
\label{qmix:main}
Let $a>0$ be algebraic and $\nu\ge0$ an integer. Every eigenvalue of
$\widehat M_{a,\nu}(H)$ algebraic over $\Kexp$ equals $1$.
In particular, $\widehat M_{a,\nu}(H)$ is nonsingular.
If the support graph of $H$, with edges $uv$ whenever $u\ne v$ and
$H_{uv}\ne0$, has $c$ connected components, eigenvalue $1$ has
multiplicity exactly $c$. Every other eigenvalue is real, belongs to
$(-1,1)$, and is transcendental over $\Kexp$.
\end{theorem}

\begin{proof}
The Poisson integral for Bessel functions \cite[\S10.9]{DLMF} gives
the characteristic function
\begin{equation}
\label{qmix:characteristic}
 \phi_{a,\nu}(s)=\Ex[e^{isR}]
 =\frac{2^\nu\nu!}{(as)^\nu}J_\nu(as)\quad(s\ne0),
 \qquad \phi_{a,\nu}(0)=1.
\end{equation}
For distinct positive algebraic $x_j$, Theorem~\ref{cf:joint}, applied
at $ix_j$, makes all pairs $(J_0(x_j),J_1(x_j))$ independent over
$\Kexp$. The recurrence
$J_{n+1}(x)=(2n/x)J_n(x)-J_{n-1}(x)$ expresses $J_\nu(x)$ as a
nonzero algebraic linear form in its pair: successive-pair recurrence
matrices have determinant $1$. Therefore, for distinct positive
algebraic $\Delta$, the values $z_\Delta=\phi_{a,\nu}(\Delta)$ are
jointly algebraically independent over $\Kexp$.

Write $H=\sum_r\theta_rE_r$ using distinct eigenvalues. All $\theta_r$
and projector entries are algebraic, since
$E_r=\prod_{s\ne r}(H-\theta_sI)/(\theta_r-\theta_s)$.
Let $D$ be the set of distinct positive gaps $|\theta_r-\theta_s|$,
and define real symmetric algebraic matrices
\begin{align*}
 B_0&=\sum_r E_r\circ\overline{E_r},\\
 B_\Delta&=\sum_{\substack{r<s\\|\theta_r-\theta_s|=\Delta}}
 \bigl(E_r\circ\overline{E_s}+E_s\circ\overline{E_r}\bigr).
\end{align*}
The characteristic-function expansion of the average, as in
\cite[Proposition~13]{BCM2024}, gives
\begin{equation}
\label{qmix:formal}
 \widehat M_{a,\nu}(H)=P(z),\qquad
 P(Z)=B_0+\sum_{\Delta\in D}Z_\Delta B_\Delta,
 \qquad P(\mathbf1)=I_N.
\end{equation}
The last identity follows from $\sum_rE_r=I_N$.

If an eigenvalue $\rho$ is algebraic over $\Kexp$, the tuple $z$
remains algebraically independent over $L=\Kexp(\rho)$.
Indeed $L/\Kexp$ and $L(z)/\Kexp(z)$ are algebraic, so
$\trdeg_L L(z)=\trdeg_{\Kexp}\Kexp(z)=|D|$.
Thus evaluation at $z$ is injective on $L[Z]$. The equality
$\det(P(z)-\rho I_N)=0$ forces $\det(P(Z)-\rho I_N)$ to vanish
identically. Substituting $Z=\mathbf1$ gives $(1-\rho)^N=0$,
so $\rho=1$. In particular zero cannot be an eigenvalue.

Unitarity makes every $M_H(t)$ nonnegative and doubly stochastic.
Moreover $M_H(-t)=M_H(t)^T$, so symmetric sampling makes the average
real symmetric even when $H$ has nonreal entries.
For a support edge, $U_H(t)_{uv}=itH_{uv}+O(t^2)$, whence the
averaged entry is positive; all diagonal entries are positive as well.
Each support component consequently gives an irreducible nonnegative
stochastic block with positive diagonal. Perron--Frobenius proves the
remaining assertions.
\end{proof}

\begin{corollary}[Entries, principal minors, and uniform mixing]
\label{qmix:principal}
Under Theorem~\ref{qmix:main}, put $A=\widehat M_{a,\nu}(H)$.
\begin{enumerate}
\item If $A_{uv}$ is algebraic over $\Kexp$, then $A_{uv}=\delta_{uv}$.
Every entry strictly between $0$ and $1$ is therefore transcendental
over $\Kexp$.
\item Every principal submatrix of $A$ is nonsingular. Its determinant
is transcendental over $\Kexp$ if its index set contains a nonisolated
support vertex; otherwise that submatrix is the identity.
\item If $N>1$, the average cannot be the uniform matrix
$N^{-1}\mathbf1\mathbf1^T$.
\end{enumerate}
\end{corollary}

\begin{proof}
More generally, if a polynomial $F$ in matrix entries has coefficients
in $\Kexp$ and $F(A)$ is algebraic over $\Kexp$, the same injectivity
argument applied to $F(P(Z))-F(A)$ gives
\begin{equation}
\label{qmix:invariant}
 F(A)=F(I_N).
\end{equation}
This proves the assertion about entries. For a principal index set $S$,
$\det P_S(\mathbf1)=1$, so $\det P_S(z)\ne0$.
To decide whether this determinant polynomial is constant, let the
half-width $b$ tend to zero while $H,\nu$ remain fixed. Set
$Q_{uv}=|H_{uv}|^2$ for $u\ne v$ and
$Q_{uu}=-\sum_{v\ne u}|H_{uv}|^2$.
Symmetry of the law and its variance $b^2/(2\nu+2)$ yield
\begin{align*}
 \widehat M_{b,\nu}(H)&=I_N+\frac{b^2}{2\nu+2}Q+O(b^4),\\
 \det\bigl(\widehat M_{b,\nu}(H)_S\bigr)
 &=1-\frac{b^2}{2\nu+2}
       \sum_{u\in S}\sum_{v\ne u}|H_{uv}|^2+O(b^4).
\end{align*}
The coefficient is strictly negative exactly when $S$ contains a
nonisolated vertex. In that case $\det P_S(Z)$ is nonconstant, and
a nonconstant polynomial in an independent tuple is transcendental
over the coefficient field. If every vertex in $S$ is isolated, the
principal submatrix is the identity. Finally, the uniform matrix has
rank one, whereas $A$ has rank $N$.
\end{proof}

\begin{remark}[The algebraic-width hypothesis]
\label{qmix:p2}
For the two-vertex path the average is
\[
 \frac12\begin{pmatrix}
 1+\phi_{a,\nu}(2)&1-\phi_{a,\nu}(2)\\
 1-\phi_{a,\nu}(2)&1+\phi_{a,\nu}(2)
 \end{pmatrix},
\]
with eigenvalues $1$ and $\phi_{a,\nu}(2)$.
At a positive zero of $J_\nu(2a)$ it has rank one and is uniform;
all such half-widths are transcendental by the independence input.
For the arcsine law the second eigenvalue can also be negative, so
the nonsingularity theorem does not follow from positive definiteness.
\end{remark}

\begin{remark}[Scope of the application]
\label{qmix:scope}
Coutinho--Guo \cite[Open Problem~4.6]{CG2024} propose developing
the spectral theory of random-time average mixing matrices.
Theorem~\ref{qmix:main} settles rank and algebraic eigenvalue
possibilities for the family \eqref{qmix:density} at algebraic widths;
it does not settle their problem for arbitrary distributions or
questions about the ordinary long-time average.
The characteristic-function framework is due to
Baptista--Coutinho--Marques \cite{BCM2024}.
For deterministic algebraic times, related exclusions over $\Qa$
already follow from Lindemann--Weierstrass. The present application
uses continuous sampling and proves the stronger exclusion over
$\Kexp$: deterministic finite-system mixing matrices have entries
in $\Kexp$ and hence eigenvalues algebraic over that field.
No assertion of joint independence of the averaged eigenvalues,
or of historical priority for this application, is intended.
\end{remark}

% Research transfer note, 22 October 2026.  Not inserted into the main paper.
% Inputs from the 114-page unified manuscript: cf:joint, cxi:main, cru:main.
% New consequences are conditional only on the correctness of those inputs.
\section{Arithmetic restrictions on Robin spectra of Euclidean balls}
\label{sp:section}

Let $k=\overline{\mathbb Q}$, let $\mathbf E$ be the ring of values of
$E$-functions at algebraic arguments, and let $\overline{\mathbf E}$ be
its integral closure in $\mathbb C$.  Write
$K_{\exp}=k(e^a:a\in k)$.  We quote two further inputs from the unified
manuscript \cite[Theorems 22.2 and 34.2]{Unified}:
\begin{equation}\label{sp:inputs}
 \overline{\mathbf E}\cap k(C(t))=k\quad(t\in k^\times),\qquad
 \overline{\mathbf E}^{\times}\cap K_{\exp}
       =\{a e^b:a\in k^\times,\ b\in k\}.
\end{equation}
Here $C(t)=I_1(t)/I_0(t)$.  The first equality is a subfield consequence
of the complex one-block intersection theorem; the second follows from
the complete Bessel-field unit theorem.  Algebraic independence of
$J_0(x),J_1(x)$ at nonzero algebraic $x$ is also available, by putting
$t=ix$ in the joint theorem.

Consider
\begin{equation}\label{sp:robin}
 -\Delta u=\lambda u\quad\hbox{in }B_R\subset\mathbb R^d,
 \qquad \partial_nu+\sigma u=0\quad\hbox{on }\partial B_R,
\end{equation}
where $d\ge2$ and $R>0$ is algebraic.  In the statements about the real
spectrum, $\sigma$ is real.  The zero eigenvalue is explicitly excluded
from transcendence assertions: it occurs when $R\sigma=-\ell$ for a
nonnegative integer angular degree $\ell$.

\subsection{An exact even-dimensional obstruction}

\begin{lemma}\label{sp:integer-ratio}
For $x\in k^\times$ and $n\in\mathbb Z_{\ge0}$, $J_n(x)\ne0$ and
\[
 L_n(x):=\frac{xJ_n'(x)}{J_n(x)}\notin\overline{\mathbf E}.
\]
In fact $k(L_n(x))=k(C(ix))$ and
$\overline{\mathbf E}\cap k(L_n(x))=k$.
\end{lemma}
\begin{proof}
The recurrence
\[
 \binom{J_{n+1}(x)}{J_{n+2}(x)}
 =\begin{pmatrix}0&1\\-1&2(n+1)/x\end{pmatrix}
       \binom{J_n(x)}{J_{n+1}(x)}
\]
has determinant one.  It expresses $(J_n,J_{n+1})$ as an invertible
$k$-linear transformation of $(J_0,J_1)$ at the fixed algebraic point.
Independence of $J_0(x),J_1(x)$ gives nonvanishing of each $J_n(x)$.
Moreover
\[
 L_n(x)=n-x\frac{J_{n+1}(x)}{J_n(x)}
\]
is a nonconstant fractional-linear transformation of
$J_1(x)/J_0(x)=C(ix)/i$.  Thus the fields agree, and
\eqref{sp:inputs} gives the result.
\end{proof}

\begin{theorem}[Even-dimensional Robin spectra]\label{sp:even}
If $d$ is even and $\sigma\in\overline{\mathbf E}$, every nonzero
eigenvalue of \eqref{sp:robin} is transcendental.
\end{theorem}
\begin{proof}
Suppose $\lambda\ne0$ is algebraic and set $x=R\sqrt\lambda\in k^\times$.
Separation into spherical harmonics gives, for some $\ell\ge0$, a regular
radial solution proportional to
$r^{-(d-2)/2}J_n(xr/R)$, with $n=\ell+(d-2)/2\in\mathbb Z_{\ge0}$.
The boundary equation is
\[
 \frac{xJ_n'(x)}{J_n(x)}=\frac{d-2}{2}-R\sigma.
\]
Its right side lies in $\overline{\mathbf E}$, contradicting
Lemma~\ref{sp:integer-ratio}.  The argument allows complex $\sigma$ and
complex $\lambda$ as well.
\end{proof}

For a fixed nonzero algebraic spectral parameter, the angular
eigenvalues of the Dirichlet-to-Neumann map on an even-dimensional ball
are
\[
 \Lambda_\ell(\lambda)=\frac1R
 \left(\frac{xJ_{\ell+(d-2)/2}'(x)}
 {J_{\ell+(d-2)/2}(x)}-\frac{d-2}{2}\right).
\]
Each generates the same field $k(C(ix))$.  Consequently the field
generated by any nonempty collection of these angular eigenvalues has
intersection exactly $k$ with $\overline{\mathbf E}$.  This is an
arithmetic assertion about the boundary response operator, stronger
than the individual transcendence of its eigenvalues.

\subsection{Positive eigenvalues in every dimension}

\begin{lemma}[Two-exponential logarithmic derivatives]
\label{sp:twoexp}
Let $x\in k^\times$, and suppose a value and a derivative have the form
\[
 f=a e^x+b e^{-x},\qquad f_1=A e^x+B e^{-x},
 \qquad a,b,A,B\in k,\quad aB-bA\ne0.
\]
If $f\ne0$ and $f_1/f\in\overline{\mathbf E}$, then $ab=0$.
\end{lemma}
\begin{proof}
Put $r=f_1/f$.  Inverting the coefficient matrix gives
\[
 \frac{e^x}{f}=\frac{B-br}{aB-bA}\in\overline{\mathbf E}.
\]
Since $e^{-x}\in\mathbf E$, $1/f\in\overline{\mathbf E}$.
Also $f\in\mathbf E$, so it is a unit of $\overline{\mathbf E}$ lying
in $K_{\exp}$.  By \eqref{sp:inputs}, $f=\alpha e^\beta$ with
$\alpha\in k^\times$, $\beta\in k$.  Lindemann--Weierstrass, applied
to the distinct exponents among $x,-x,\beta$, then forces $ab=0$.
\end{proof}

\begin{theorem}[All positive Robin eigenvalues]\label{sp:positive}
For every $d\ge2$ and every real
$\sigma\in\overline{\mathbf E}$, all positive eigenvalues of
\eqref{sp:robin} are transcendental.
\end{theorem}
\begin{proof}
Only odd $d=2q+3$ remains.  For angular degree $\ell$, put
$N=\ell+q$.  The regular radial solution for a positive eigenvalue is
$r^{-q}j_N(xr/R)$, where $x=R\sqrt\lambda$ and $j_N$ is the spherical
Bessel function.  Thus
\[
 R\sigma=q-\frac{xj_N'(x)}{j_N(x)}.
\]
At real algebraic $x>0$, the finite Hankel expansions give
$j_N(x)=a e^{ix}+b e^{-ix}$, and its derivative has the same two
exponents with algebraic coefficients.  The determinant of these two
coefficient rows is nonzero by the Hankel Wronskian.  Moreover both
$a$ and $b$ are nonzero: the outgoing and incoming Hankel functions are
complex conjugates on the positive real axis, and a common zero would
contradict their nonzero Wronskian.  The boundary equation would put
$j_N'(x)/j_N(x)$ in $\overline{\mathbf E}$, contradicting
Lemma~\ref{sp:twoexp}, with $ix$ in place of $x$.
\end{proof}

\subsection{A complete description of the odd-dimensional exceptions}

For $N\ge0$ define the explicit polynomials
\begin{equation}\label{sp:polynomials}
 a_{N,k}=\frac{(N+k)!}{2^k k!(N-k)!},\qquad
 P_N(y)=\sum_{k=0}^N(-1)^k a_{N,k}y^{N-k},\qquad
 Q_N(y)=\sum_{k=0}^N a_{N,k}y^{N-k}.
\end{equation}
The modified spherical Bessel function is
\begin{equation}\label{sp:ispherical}
 i_N(y)=\frac{e^yP_N(y)+(-1)^{N+1}e^{-y}Q_N(y)}{2y^{N+1}}.
\end{equation}
These are the terminating half-integer Hankel formulas; see DLMF
\S\S10.47--10.49.

\begin{theorem}[Exact negative algebraic exceptions]\label{sp:negative}
Let $d=2q+3$ and $\sigma\in\overline{\mathbf E}\cap\mathbb R$.
An algebraic number $\lambda<0$ is an eigenvalue of
\eqref{sp:robin} if and only if, for some integer $N\ge q$, the number
$y=R\sqrt{-\lambda}>0$ satisfies
\begin{equation}\label{sp:exception}
 P_N(y)=0,\qquad
 R\sigma=q-y\frac{i_N'(y)}{i_N(y)}.
\end{equation}
Every such $N$ is odd, every such $y$ is algebraic, and the corresponding
Robin constant belongs to $\mathbf E$.  More explicitly,
\begin{equation}\label{sp:explicit-sigma}
 R\sigma=q+y+N+1-y\frac{Q_N'(y)}{Q_N(y)}
                   -y e^{2y}\frac{P_N'(y)}{Q_N(y)}.
\end{equation}
The coefficient of $e^{2y}$ in this formula is nonzero, so the Robin
constant of each exception is transcendental.
\end{theorem}
\begin{proof}
The regular radial solution is $r^{-q}i_N(yr/R)$ with $N=\ell+q$.
The boundary equation is the second equality in \eqref{sp:exception}.
If $\lambda$ is algebraic, then $y$ is positive algebraic, and this
equation puts $i_N'(y)/i_N(y)$ in $\overline{\mathbf E}$.
The two exponential components in \eqref{sp:ispherical} are independent
solutions of the modified spherical Bessel equation.  Their Wronskian
is nonzero at $y\ne0$; hence the coefficient determinant for
$i_N,i_N'$ is nonzero.  Lemma~\ref{sp:twoexp} now implies that one
exponential coefficient vanishes.  Since $Q_N(y)>0$, it must be
$P_N(y)$.

Conversely, a positive root of $P_N$ is algebraic, and
\eqref{sp:ispherical} is nonzero there because $Q_N(y)>0$.
The boundary constant defined in \eqref{sp:exception} therefore gives
an eigenfunction of angular degree $N-q$.  Positivity $i_N(y)>0$,
which follows from the series for $I_{N+1/2}$, implies that $N$ is odd
when $P_N(y)=0$.  Differentiating \eqref{sp:ispherical} at such a root
gives \eqref{sp:explicit-sigma}.  Its coefficients are algebraic, so
$\sigma\in k[e^{2y}]\subset\mathbf E$.

Finally $P_N'(y)\ne0$.  Otherwise the component
$e^yP_N(y)/(2y^{N+1})$ and its derivative would both vanish at the
ordinary point $y>0$ of its second-order differential equation.
Uniqueness would make that component identically zero, a contradiction.
Lindemann--Weierstrass consequently makes the displayed boundary
constant transcendental.
\end{proof}

\begin{example}[A necessary exception, not a numerical artifact]
In the unit ball in $\mathbb R^3$, take $N=1$.  Then
$P_1(y)=y-1$, $Q_1(y)=y+1$, and
\[
 \boxed{\quad \sigma=\frac{5-e^2}{2}\in\mathbf E,\qquad
                   \lambda=-1.\quad}
\]
Indeed $i_1(1)=e^{-1}$ and
$i_1'(1)=(e-5e^{-1})/2$.  The angular-degree-one eigenfunctions have
the usual three-dimensional spherical-harmonic multiplicity.
Thus the positive-spectrum theorem cannot be extended to all nonzero
eigenvalues in odd dimensions.
\end{example}

\subsection{Algebraic Robin parameters and accidental degeneracy}

The next statement uses classical transcendence and a recurrence
calculation; unlike Theorems~\ref{sp:even}--\ref{sp:negative}, its proof
does not need the new integral-intersection or unit results.

\begin{theorem}[An algebraic-parameter Bourget statement]
\label{sp:bourget}
Let $\nu\in\{0,\tfrac12\}$, $h\in k$, and let $n,m$ be distinct
nonnegative integers.  The functions
\[
 D_n(z)=zJ_{\nu+n}'(z)+hJ_{\nu+n}(z),\qquad
 D_m(z)=zJ_{\nu+m}'(z)+hJ_{\nu+m}(z)
\]
have no common nonzero complex zero.
\end{theorem}
\begin{proof}
Write $J_{\nu+n}=A_nJ_\nu+B_nJ_{\nu+1}$ with
$A_n,B_n\in\mathbb Q(z)$, using the usual recurrence.
In particular these rational functions are Laurent polynomials with
their only possible finite pole at zero.  The transfer matrices have
determinant one.  Write
$D_n=a_nJ_\nu+b_nJ_{\nu+1}$.
We first check that
\[
 \Delta_{n,m}(z)=a_n(z)b_m(z)-b_n(z)a_m(z)
\]
is not identically zero for any $h$.

Induction in the recurrence gives the following Laurent expansions at
infinity.  If $n$ is even and $\epsilon_n=(-1)^{n/2}$, then
\[
 (a_n,b_n)=\epsilon_n\left(
 h+\nu-\nu n-\frac{n^2}{2}+O(z^{-2}),\quad
 -z+O(z^{-1})\right).
\]
If $n$ is odd and $\epsilon_n=(-1)^{(n-1)/2}$, then
\[
 (a_n,b_n)=\epsilon_n\left(
 z+O(z^{-1}),\quad h-\nu n-\frac{n^2+1}{2}+O(z^{-2})\right).
\]
For clarity, the induction follows from
\[
 A_{n+1}=\frac{2(\nu+n)}z A_n-A_{n-1},\qquad
 B_{n+1}=\frac{2(\nu+n)}z B_n-B_{n-1},
\]
starting with $(A_0,B_0)=(1,0)$ and $(A_1,B_1)=(0,1)$, and then using
$(a_n,b_n)=(\nu+n+h)(A_n,B_n)-z(A_{n+1},B_{n+1})$.
Consequently, for $n,m$ of the same parity,
\[
 \Delta_{n,m}(z)=\epsilon_n\epsilon_m
 \frac{(\nu+n)^2-(\nu+m)^2}{2}\,z+O(z^{-1}),
\]
whose leading coefficient is nonzero.  For opposite parities its
leading term is $\pm\epsilon_n\epsilon_m z^2$.
Thus $\Delta_{n,m}\ne0$, independently of $h$.

If $z_0\ne0$ were a common zero, the nonzero vector
$(J_\nu(z_0),J_{\nu+1}(z_0))$ would lie in the kernel of the coefficient
matrix.  The vector cannot be zero, by the recurrence and uniqueness
for Bessel's differential equation.  Therefore
$\Delta_{n,m}(z_0)=0$, so $z_0$ is algebraic when $h$ is algebraic.

For $\nu=0$, independence of $J_0(z_0),J_1(z_0)$ contradicts
$D_n(z_0)=0$: the row $(a_n(z_0),b_n(z_0))$ is nonzero since the
transfer matrix is invertible and $z_0\ne0$.
For $\nu=\tfrac12$, use the elementary sine and cosine expressions
for $J_{1/2},J_{3/2}$.  Their coefficient matrix in
$e^{iz_0},e^{-iz_0}$ is invertible, up to a common nonzero factor.
Thus $D_n(z_0)=0$ would be a nontrivial algebraic linear relation
between $e^{iz_0}$ and $e^{-iz_0}$, contradicting
Lindemann--Weierstrass.
\end{proof}

\begin{corollary}[Removal of genericity for algebraic Robin constants]
\label{sp:pleijel}
For a Euclidean ball of algebraic radius and every real algebraic
Robin constant, there are no accidental coincidences between distinct
spherical-harmonic degrees.  Hence every eigenfunction is a separated
radial factor times a spherical harmonic, and the Robin Pleijel
constant equals the Dirichlet Pleijel constant by the theorem of
Hassannezhad--Sher.  In dimension two this common constant is
approximately $0.461$.
\end{corollary}
\begin{proof}
The radial boundary condition has the form in
Theorem~\ref{sp:bourget} with
$h=R\sigma-(d-2)/2$ and with integer or half-integer orders.
The common-zero theorem excludes all nonzero coincidences.  At zero
energy the boundary condition on a harmonic homogeneous polynomial of
degree $\ell$ is $R\sigma=-\ell$, so distinct degrees cannot coincide
there either.  Apply Hassannezhad--Sher, Theorem~1.4.
\end{proof}

\subsection{Known status and the actual remaining obstruction}

Hassannezhad--Sher~\cite{HS2024}, Assumption~1.3 and Theorem~1.4, prove the
ball Pleijel equality under absence of accidental coincidences.  Their
Proposition~3.2 makes that condition generic, and Remark~3.3 gives a
real Robin parameter for which it fails.  Their Open Question~1.5
asks for boundary-condition independence on general domains.
Corollary~\ref{sp:pleijel} supplies an explicit arithmetic subclass;
it does not settle that general-domain question.  Wang~\cite{Wang2023},
Remark~1.3, explicitly records the disk noncoincidence result for
algebraic Robin constants less than one.  The calculation above has
no such restriction.  A literature search has not established priority
for this unrestricted elementary extension, so it should not be
presented as the solution of a newly discovered named conjecture.

The clamped round-plate problem is different.  Its characteristic
functions are
$W_n(x)=I_n'(x)J_n(x)-I_n(x)J_n'(x)$, and the desired multiplicity-two
bound would rule out a common positive zero of two distinct $W_n$.
Mangoubi--Rosenblatt~\cite{MR2026}, Conjecture~1.3, leaves precisely
this problem open and proves multiplicity at most four in dimension
two.  A common zero is already transcendental.  The unified paper's
arithmetic value theorems apply at algebraic arguments and therefore
cannot simply be evaluated at that zero.  Their second-order
recurrence produces a polynomial relation between the common zero and
a Bessel quotient at that zero, not a polynomial over $k$ in the zero
alone.  No removal of this obstruction is proved here.

Likewise, for an $\overline{\mathbf E}$-valued Robin parameter, the
recurrence determinant makes a hypothetical common wavenumber integral
over $\overline{\mathbf E}$ after clearing a power of $z$: its leading
coefficient is a nonzero algebraic constant.  It need not make that
wavenumber algebraic.  Applying an algebraic-argument value theorem
at such a number would be invalid.  Thus the no-accidental-degeneracy
statement above is asserted for algebraic parameters, not for every
parameter in $\overline{\mathbf E}$.


\begin{thebibliography}{99}

\bibitem{Unified}
Keio, \emph{Euler--Gompertz and Bessel Values: Joint Independence,
Integral Intersections, and Critical Limits}, revised integrated research
manuscript, 22 October 2026, 114 pp. Working source supplied with this project,
\texttt{Keio\_Euler\_Bessel\_Unified\_2026.tex}; not a published reference.

\bibitem{CG2024}
G. Coutinho and K. Guo, \emph{Selected open problems in continuous-time
quantum walks}, Special Matrices \textbf{12} (2024), 20240025.
\href{https://doi.org/10.1515/spma-2024-0025}{doi:10.1515/spma-2024-0025}.

\bibitem{BCM2024}
P. Baptista, G. Coutinho and V. Marques,
\emph{Unexpected averages of mixing matrices}, Quantum Information
Processing \textbf{23} (2024), 272.
\href{https://arxiv.org/abs/2308.16378}{arXiv:2308.16378}.

\bibitem{DLMF}
NIST Digital Library of Mathematical Functions, \S10.6, \S10.9,
\S10.47, \S10.49, \url{https://dlmf.nist.gov/10}.

\bibitem{HS2024}
A. Hassannezhad and D. Sher, \emph{On Pleijel's nodal domain theorem
for the Robin problem}, Bulletin of the London Mathematical Society
\textbf{56} (2024), 1449--1467.
\href{https://doi.org/10.1112/blms.13005}{doi:10.1112/blms.13005}.

\bibitem{Wang2023}
W. Wang, \emph{A uniform bound on the multiplicity of Dirichlet
eigenvalues in planar annuli}, Journal of Mathematical Analysis and
Applications \textbf{518} (2023), 126616.
\href{https://doi.org/10.1016/j.jmaa.2022.126616}{doi:10.1016/j.jmaa.2022.126616}.

\bibitem{MR2026}
D. Mangoubi and D. Rosenblatt, \emph{On multiplicity bounds for eigenvalues
of the clamped round plate}, Journal of Spectral Theory \textbf{16}
(2026), 183--195.
\href{https://doi.org/10.4171/JST/588}{doi:10.4171/JST/588}.

\bibitem{FHW}
M.~D. Fraser, Y.-S. Hsu and J.~J. Walker,
\emph{Identifiability of finite mixtures of von Mises distributions},
Annals of Statistics \textbf{9} (1981), 1130--1131.
\href{https://doi.org/10.1214/aos/1176345595}{doi:10.1214/aos/1176345595}.

\bibitem{Kent}
J.~T. Kent, \emph{Identifiability of finite mixtures for directional data},
Annals of Statistics \textbf{11} (1983), 984--988.
\href{https://doi.org/10.1214/aos/1176346264}{doi:10.1214/aos/1176346264}.

\bibitem{HG}
K. Hornik and B. Gr\"un,
\emph{On maximum likelihood estimation of the concentration parameter
of von Mises--Fisher distributions}, Computational Statistics
\textbf{29} (2014), 945--957.
\href{https://doi.org/10.1007/s00180-013-0471-0}{doi:10.1007/s00180-013-0471-0}.

\bibitem{YX}
Z. Yang and L. Xie, \emph{Frequency-selective Vandermonde decomposition
of Toeplitz matrices with applications}, version 3, 2017.
\href{https://arxiv.org/abs/1605.02431}{arXiv:1605.02431}.
The classical positive Toeplitz mechanism is recalled there; our
finite-atomic argument is proved directly above.

\bibitem{ADS}
C. Am\'endola, M. Drton and B. Sturmfels,
\emph{Maximum likelihood estimates for Gaussian mixtures are
transcendental}, version 2, 2015.
\href{https://arxiv.org/abs/1508.06958}{arXiv:1508.06958}.

\bibitem{Ng}
T.~L.~J. Ng, \emph{Penalized maximum likelihood estimator for mixture
of von Mises--Fisher distributions}, version 2, 2022.
\href{https://arxiv.org/abs/2009.02921}{arXiv:2009.02921}.

\end{thebibliography}
\end{document}
